%% kapa: exam=1 round=19 year=2008 period=1 form=A subject=civil_law subject_ko=민법 seq=40 number=40
%% source: Q-Net 감정평가사 기출문제 (cst00309) artlSeq=1001962 file=extracted/1008/1001962/2008년 제19회 감정평가사 1차시험 기출문제/1교시.hwp
%% conversion: kordoc 4.24.51 -> tools/gen1.py (auto); review=proofread-sonnet; flags=none
\begin{question}{40}
가등기담보 등에 관한 법률(이하 ‘가등기담보법’이라 함)에 관한 설명으로 틀린 것은? (다툼이 있는 경우에는 판례에 의함)
\begin{choices}
  \item 공사대금채권을 담보할 목적으로 가등기가 경료된 경우에는 가등기담보법이 적용되지 않는다.
  \item 가등기담보법은 등기 또는 등록할 수 있는 부동산소유권 이외의 권리의 취득을 목적으로 하는 담보계약에도 준용되지만, 저당권의 취득을 목적으로 하는 담보계약에는 준용되지 않는다.
  \item 가등기담보권자가 채무자 등에게 통지한 청산금의 액수가 객관적인 청산금의 평가액에 미치지 못하더라도, 담보권실행의 통지로서의 효력이나 청산기간의 진행에는 아무런 영향이 없다.
  \item 가등기담보권자의 청산금지급채무와 가등기담보권설정자의 소유권이전등기 및 목적물인도채무는 동시이행의 관계에 있다.
  \item 채권자의 청산금평가액 자체에 이의가 있는 후순위권리자는 청산기간 내에 자기 채권의 변제기가 도래한 경우에 한하여 독자적으로 경매를 청구할 수 있다.
\end{choices}
\end{question}
\answer{5}
